% Figure for Oscillations, Waves and Optics §A6.1 (double-slit geometry: (a) slits S1, S2 a distance d apart, screen at distance D, point P at height y and angle theta; (b) close-up: for a distant screen the rays are parallel and differ in length by d sin(theta)).
\begin{tikzpicture}[font=\small, >={Stealth[length=2.2mm]}]
  % ---------- panel (a): overall layout (not to scale) ----------
  \begin{scope}
    \fill[black!30] (-0.12,-2.0) rectangle (0.12,-0.62);
    \fill[black!30] (-0.12,-0.38) rectangle (0.12,0.38);
    \fill[black!30] (-0.12,0.62) rectangle (0.12,2.0);
    \node[font=\scriptsize, left] at (-0.12,0.5) {$S_1$};
    \node[font=\scriptsize, left] at (-0.12,-0.5) {$S_2$};
    \draw[black!55, thin] (0,0) -- (5.2,0);
    \draw[very thick, black!70] (5.2,-2.0) -- (5.2,2.0) node[above, font=\scriptsize] {screen};
    \coordinate (P) at (5.2,1.55);
    \draw[figB, thick] (0.12,0.5) -- (P) node[pos=0.55, above, sloped, font=\scriptsize] {$r_1$};
    \draw[figB, thick] (0.12,-0.5) -- (P) node[pos=0.6, below, sloped, font=\scriptsize] {$r_2$};
    \draw[black!55, dashed, thin] (0,0) -- (P);
    \fill (P) circle (1.6pt) node[above right, font=\scriptsize] {$P$};
    \draw[black!55, thin] (1.3,0) arc (0:{atan(1.55/5.2)}:1.3);
    \node[font=\scriptsize] at (1.55,0.2) {$\theta$};
    \draw[<->, black!60, thin] (5.5,0) -- (5.5,1.55) node[midway, right, font=\scriptsize] {$y$};
    \draw[<->, black!60, thin] (0,-2.3) -- (5.2,-2.3) node[midway, below, font=\scriptsize] {$D$};
    \draw[<->, black!60, thin] (-0.75,-0.5) -- (-0.75,0.5) node[midway, left, font=\scriptsize] {$d$};
    \node[font=\scriptsize, black!70] at (2.6,-3.0) {(a) layout, $D \gg d$ (not to scale)};
  \end{scope}
  % ---------- panel (b): close-up near the slits ----------
  \begin{scope}[xshift=8.3cm]
    \def\th{25}
    \fill[black!30] (-0.12,-2.0) rectangle (0.12,-1.0);
    \fill[black!30] (-0.12,-0.75) rectangle (0.12,0.75);
    \fill[black!30] (-0.12,1.0) rectangle (0.12,2.0);
    \coordinate (A) at (0.12,0.875);
    \coordinate (B) at (0.12,-0.875);
    \node[font=\scriptsize, left] at (-0.12,0.875) {$S_1$};
    \node[font=\scriptsize, left] at (-0.12,-0.875) {$S_2$};
    \draw[figB, thick, ->] (A) -- ++(\th:4.2) node[right, font=\scriptsize] {to $P$};
    \draw[figB, thick, ->] (B) -- ++(\th:5.0) node[right, font=\scriptsize] {to $P$};
    % foot of the perpendicular from S1 onto ray 2
    \coordinate (F) at ($(B)+(\th:{1.75*sin(\th)})$);
    \draw[black!55, dashed, thin] (A) -- (F);
    \draw[figR, line width=1.6pt] (B) -- (F);
    \node[figR, font=\scriptsize, anchor=north west] at ($(B)+(0.3,-0.12)$) {$\Delta l = d\sin\theta$};
    % right-angle mark at F
    \draw[black!55, thin] ($(F)+(\th:-0.18)$) -- ++(\th+90:0.18) -- ++(\th:0.18);
    % angle theta at S2 between the axis and the ray
    \draw[black!55, thin] (B) -- ++(1.6,0);
    \draw[black!55, thin] ($(B)+(1.25,0)$) arc (0:\th:1.25);
    \node[font=\scriptsize] at ($(B)+(12:1.5)$) {$\theta$};
    % angle theta at S1 between the barrier and the perpendicular
    \draw[black!55, thin] ($(A)+(0,-0.55)$) arc (-90:{-90+\th}:0.55);
    \node[font=\scriptsize] at ($(A)+({-90+\th/2}:0.8)$) {$\theta$};
    \draw[<->, black!60, thin] (-1.05,-0.875) -- (-1.05,0.875) node[midway, left, font=\scriptsize] {$d$};
    \node[font=\scriptsize, black!70] at (2.2,-3.0) {(b) close-up: rays to a distant $P$ are parallel};
  \end{scope}
\end{tikzpicture}
